\scanpage{46}\mspage{45}

axiomes qui vont suivre, il y en a qui vont éliminer, je pense, cet exemple trivial d'atelier.

\par\noindent\hfill\rule{1.4cm}{.3pt}\hfill\mbox{}\par

Revenons aux \uline{petites} ombres et multiombres. C'est le moment d'introduire enfin

\uline{At 8} Soit $X\in M$, et soit $\unsure{\delta|\partial} X=\mathop{\mathrm{Sup}_{\deleted{\unsure{\mathrm{au}|\mathrm{aux}}}}}\limits_{\substack{Y\mathrel{\unsure{\triangleleft|\trianglelefteq}}X\\Y\ne X}}^{(\leq)}\,Y$.

Alors $\exists\,x\in\unsure{\mathcal L|L}$ tel que $x\unsure{\ll|\leq}X$, \deleted{$x\unsure{\in|<}$} $x\mathrel{\not\ll}\delta X$ \deleted{\unsure{et}} (i.e. $x\mathrel{\overset{\circ}{\ll}}X$). \deleted{et} \deleted{\unsure{ou un} En}

En d'autres termes, posant \deleted{\unsure{p}}
\[
\begin{aligned}
\mathrm{omb}(X)^\circ&=\mathrm{omb}(X)-\mathrm{omb}(\delta X)\\
&=\mathrm{omb}(X)-\bigcup_{\substack{Y\mathrel{\triangleleft} X\\Y\ne X}}\mathrm{omb}(Y)\unsure{,}\ ,
\end{aligned}
\]
on a :
\[
\boxed{\mathrm{omb}(X)^\circ\ne\varnothing}
\tag*{\makebox[0pt][r]{\makebox[\displaywidth][l]{(1.52)}}}
\]

Considérons alors, pour $F\in\mathcal F$, \deleted{\unsure{sa}} multiombre, qui est par definition \deleted{formée des \unsure{part} parties} \inserted{la figure \unsure{ens.iste|ensembliste}}
\[
\mathrm{mulomb}(F)=\{\mathrm{omb}(X)\mid\underset{\inserted{\text{i.e. }X\in\widetilde F}}{X\mathrel{\triangleleft} F},\ \mathrm{omb}(X)^\circ\ne\varnothing\}.
\]
At 8 nous dit que l'ens d'indices à prendre ici est $\widetilde F$ tout entier. Ainsi on trouve \inserted{(cf \unsure{sorite|sorte} ensembliste)}

\uline{Corollaire} \deleted{\unsure{Question|Soient}} $F$ une figure, $X,Y\in\widetilde F$. Alors conditions suivantes équivalentes
\[
\unsure{\text{a)}|\text{.)}}\quad X\leq Y\qquad(\Longleftrightarrow X\ll Y)
\]
